\centerline{\bf \Large  7 класс}

\problem{\textbf{1}} Пюрвя назвал натуральное число \textit{УДЕкрасивым}, если сумма цифр числа равна 16, при этом никакая цифра не встречается дважды. Сколько существует четырёхзначных \textit{УДЕкрасивых} чисел, меньших 2025?

\textbf{Решение:} Т.к. речь идёт о четырёхзначных числах, то начнём перебор с 1000 и закончим 2025.
Выпишем все возможные варианты УДЕкрасивых чисел:

1069, 1078 ,1087, 1096

1249, 1258, 1267, 1276, 1285, 1294

1348, 1357, 1375, 1384

1429, 1438, 1456, 1465, 1483, 1492

1528, 1537, 1546, 1564, 1573, 1582

1609, 1627, 1645, 1654, 1672, 1690

1708, 1726, 1735, 1753, 1762, 1780

1807, 1825, 1834, 1843, 1852, 1870

1906, 1924, 1942, 1960

Всего 48 чисел.

\textbf{Критерии.} \\
\textit{7 баллов} $-$ 47-48 правильных чисел\\
\textit{6 баллов} $-$ 42-46 правильных чисел\\
\textit{5 баллов} $-$ 35-41 правильных чисел\\
\textit{4 балла} $-$ 28-34 правильных чисел\\
\textit{3 балла} $-$ 21-27 правильных чисел\\
\textit{2 балла} $-$ 14-20 правильных чисел\\
\textit{1 балл} $-$ 7-13 правильных чисел\\
\textit{0 баллов} $-$ 0-6 правильных чисел\\
\textit{0-7 баллов} $-$ Комбинаторный способ\\

\problem{\textbf{2}} Решите задачу:

а) В малом сборнике удивительных задач Пюрви Мучкаевича 4 раздела, каждый из которых содержит одинаковое количество задач. Очир правильно решил 20 заданий. При этом процент правильно решённых заданий оказался больше 60, но меньше 70. Сколько всего заданий в малом сборнике Пюрви Мучкаевича?

б) Составьте и решите обратную задачу по схеме: 4; ?; 60; 70; 32.

\textbf{Решение А:} Пусть в сборнике $x$ задач. По условию $\dfrac{60}{100}<\dfrac{20}{x}<\dfrac{70}{100}$, отсюда $28 \dfrac{4}{7}=\dfrac{200}{7}<x<\dfrac{100}{3}=33 \dfrac{1}{3}$, то есть $29 \leq x \leq 33$. Из первого условия задачи следует, что число заданий должно делиться на 4. Значит, всего 32 задания в малом сборнике Пюрви Мучкаевича.

\textbf{Примерное условие Б:} В малом сборнике удивительных задач Пюрви Мучкаевича 4 раздела, каждый из которых содержит одинаковое количество задач. Очир правильно решил несколько заданий. При этом процент правильно решённых заданий оказался больше 60, но меньше 70. Сколько всего заданий решил Очир, если в малом сборнике Пюрви Мучкаевича 32 задачи?

\textbf{Решение Б:} Пусть Очир решил $x$ задач. По условию $\dfrac{60}{100}<\dfrac{x}{32}<\dfrac{70}{100}$, отсюда $19 \dfrac{2}{10}=\dfrac{192}{10}<x<\dfrac{224}{10}=22 \dfrac{4}{10}$, то есть $20 \leq x \leq 22$. Значит всего Очир мог решить 20, 21, 22 заданий. Можно в обратной задаче добавить дополнительное условие: Очир решил одинаковое число задач из каждого раздела. Тогда это даст условие делимости на 4, количества решённых Очиром заданий. И тогда ответ будет единственным 20. В противном случае 3 ответа: 20, 21, 22.

\textbf{Критерии.} \\
\textit{0-4 баллов} $-$ Решение пункта А\\
\textit{0-1 балла} $-$ Составление условия пункта Б \\
\textit{0-2 баллов} $-$ Решение пункта Б\\
\textit{0 баллов} $-$ Любой ответ без объяснений\\

\problem{\textbf{3}} Во дворе у Батыра есть три закрытых загона с номерами 1, 2 и 3. В одном из них 1 коричневый, 1 чёрный и 1 белый ягнёнок, в другом - 1 коричневый и 1 белый, в третьем - 1 чёрный ягнёнок. При этом известно, что номер загона не совпадает с числом ягнят внутри загона. Как, узнав цвет только одного ягнёнка из одного загона, найти число ягнят в каждом загоне?

\textbf{Решение:} Узнаем цвет ягнёнка из 3 загона. В 3 загоне может быть 2 ягнёнка или 1, при этом сколько на самом деле ягнят определяется однозначно благодаря различным цветам (поэтому не узнаем цвет из других загонов). Если в 3 загоне 2 ягнёнка, то в 1 уже не может быть 2, а это значит, что в 1 загоне - 3 ягнёнка и во 2 загоне - 1 ягнёнок. Если же в 3 загоне 1 ягнёнок, то во 2 уже не может быть 1, а это значит, что во 2 загоне - 3 ягнёнка и в 1 загоне - 2 ягнёнка. 

\textbf{Критерии.} \\
\textit{1 балл} $-$ В первом 2 или 3, во втором 1 или 3 и в третьем 1 или 2 ягнёнка\\
\textit{0-2 балла} $-$ Однозначное определение последующих загонов, если знаешь любой один\\
\textit{0-4 балла} $-$ Подходить нужно к третьему (однозначное определение по цветам)\\
\textit{0-7 баллов} $-$ Полный перебор\\

\problem{\textbf{4}} В стране УДЕ мальчики днём всегда говорят правду, а ночью всегда лгут, а девочки — наоборот. Днём каждый из них сказал: “У меня знакомых мальчиков на 1 больше, чем знакомых девочек”, а ночью — “Среди незнакомых мне жителей страны мальчиков на 1 больше, чем девочек”. Могло ли в стране УДЕ быть ровно 2025 жителей?

\textbf{Решение:} Для решения задачи полезно знать следующий факт, который мы для удобства сформулируем в терминах знакомства. Не существует компании, состоящей из нечётного числа людей, в которой каждый человек имеет нечётное количество знакомых (среди людей из этой компании).

Доказательство. Пусть все знакомые пожмут друг другу руки. Так как в каждом рукопожатии участвуют две руки, значит, суммарное количество «протянутых рук» четно. С другой стороны, каждый человек протягивал руку нечётное число раз, чтобы суммарное число «протянутых рук» в этой ситуации получилось чётным, число людей должно быть чётным.

Каждый мальчик в днём говорил правду, значит, у него нечётное число знакомых (поскольку среди них мальчиков на 1 больше, чем девочек). А каждая девочка сказала правду ночью, и по той же причине у неё нечётное число незнакомых.

Предположим, что в стране 2025 жителей. Тогда у каждого человека сумма количества знакомых и незнакомых равна 2024 (чётное число). Следовательно, у каждого человека количество знакомых имеет ту же чётность, что и количество незнакомых.

Таким образом, у каждой девочки нечётное количество незнакомых и, следовательно, нечётное количество знакомых. В итоге получаем, что каждый из 2025 жителей, имеет нечётное число знакомых, чего быть не может.

\textbf{Критерии.} \\
\textit{0-5 баллов} $-$ Продвижение через идеи чётности\\
\textit{1 балл} $-$ Лемма о рукопожатиях без доказательства\\
\textit{1 балл} $-$ Доказательство леммы о рукопожатиях\\
\textit{0 баллов} $-$ Любой ответ без объяснений\\

\problem{\textbf{5}} Арслан написал 2025-значное число. Потом всеми возможными способами выбрал пару цифр, сложил цифры в каждой паре и полученные 2049300 чисел перемножил. Мог ли Арслан в результате получить исходное число? 

\textbf{Решение:} Пусть число Арслана равно $X$. Тогда, так как Арслан записал 2025-значное число, имеем $X<10^{2025}$. Заметим, что в числе Арслана не может быть больше одного нуля, иначе хотя бы одна из сумм равнялась нулю, а значит, всё произведение равнялось бы нулю. Следовательно, нулей не больше одного. Найдём минимальное искомое произведение. Оно достигается, если число состоит из 2024 единиц и одного нуля. В этом случае 2024 суммы будут равны 1, а остальные - 2. Получается, что минимальное произведение равно $1^{2024}\cdot 2^{2047276}$. Имеем следующее неравенство: $2^{2047276}\leq X<10^{2025}<16^{2025}=2^{8100}$. Чего не может быть.

\textbf{Критерии.} \\
\textit{7 баллов} $-$ Полностью верное решение\\
\textit{0-6 баллов} $-$ Продвижения в решении\\
\textit{0 баллов} $-$ Любой ответ без объяснений\\